% The two layouts, byte by byte.
% Millimetres throughout: \framefield draws at y = 0 in millimetres.
\begin{tikzpicture}[font=\sffamily\scriptsize]

% ================================================================ the state frame
\begin{scope}
  \node[signame] at (-2mm,4mm) {state, node to master};
  \framefield{0}{16}{fdata}{position\\{\tiny int32, counts}}
  \framefield{16}{16}{fdata}{velocity\\{\tiny int32, mrad/s}}
  \framefield{32}{16}{fcrc}{effort\\{\tiny int32, mNm, DERIVED}}
  \framefield{48}{16}{fid}{timestamp\\{\tiny uint32, us, wraps}}
  \framefield{64}{8}{fstuff}{valid\\{\tiny u16}}
  \framefield{72}{8}{fstuff}{faults\\{\tiny u16}}
  \framefield{80}{4}{fid}{md}
  \framefield{84}{4}{fcrc}{sr}
  \framefield{88}{4}{fcrc}{mv}
  \framefield{92}{4}{fid}{sq}
  \node[fbyte] at (8mm,11mm) {0};
  \node[fbyte] at (24mm,11mm) {4};
  \node[fbyte] at (40mm,11mm) {8};
  \node[fbyte] at (56mm,11mm) {12};
  \node[fbyte] at (68mm,11mm) {16};
  \node[fbyte] at (76mm,11mm) {18};
  \node[fbyte] at (86mm,11mm) {20 .. 23};
  \node[chlabel, anchor=west] at (98mm,4mm) {24 bytes exactly, which is};
  \node[chlabel, anchor=west] at (98mm,0mm) {an allowed length};
\end{scope}

% ================================================================ the command frame
\begin{scope}[yshift=-24mm]
  \node[signame] at (-2mm,4mm) {command, master to node};
  \framefield{0}{16}{fdata}{target\\{\tiny int32}}
  \framefield{16}{16}{fid}{valid until\\{\tiny uint32, us}}
  \framefield{32}{8}{fstuff}{mode req}
  \framefield{40}{8}{fstuff}{policy}
  \framefield{48}{8}{fcrc}{seq}
  \framefield{56}{8}{fstuff}{spare}
  \node[fbyte] at (8mm,11mm) {0};
  \node[fbyte] at (24mm,11mm) {4};
  \node[fbyte] at (36mm,11mm) {8};
  \node[fbyte] at (44mm,11mm) {9};
  \node[fbyte] at (52mm,11mm) {10};
  \node[fbyte] at (60mm,11mm) {11};
  \node[chlabel, anchor=west] at (98mm,4mm) {16 bytes, also allowed};
  \draw[busstub] (24mm,-1mm) -- (24mm,-8mm);
  \node[tlabel, anchor=north, text width=46mm, align=center] at (24mm,-8mm)
    {four bytes that change the system's failure mode: after this time the node
     holds no setpoint};
\end{scope}

% ================================================================ the rules
\node[umlnote, text width=56mm, anchor=north west] at (0mm,-48mm)
  {\textbf{Three things on the wire that are not the value.} The validity word
   says which fields are meaningful right now. The source byte says whether the
   effort figure was measured or derived, which chapter 8 insisted on. The model
   version says which parameter set produced it. Together they cost four bytes
   and stop a derived number from being read as a measurement.};

\node[umlnote, text width=56mm, anchor=north west] at (62mm,-48mm)
  {\textbf{Units chosen for exactness, not convenience.} Position travels in
   encoder counts because that is what was measured and it converts without
   loss. Velocity and effort travel in milli-units as integers. The master does
   the conversion to radians and newton metres, because the master is the one
   with the floating point and the screen.};
\end{tikzpicture}
